% 0649(본책 95쪽) — 세 집합 A·B·C 의 벤 다이어그램 6개(좌변·우변 각 3개, 결합법칙 (A∩B)∩C=A∩(B∩C)). 스캔 그림이라 TikZ 로 다시 그림.
% 원문대로 각 행의 앞 두 그림에만 색칠((A∩B)·C / A·(B∩C)), 세 번째(결과) 그림은 빈칸. 답(결과 색칠)은 적지 않는다.
\begin{tikzpicture}[scale=0.6, line width=0.5pt, every node/.style={font=\small}]
  \node[anchor=west] at (-3.7,1.7) {[좌변]};
  \begin{scope}[shift={(0,0)}]
    \begin{scope}\clip (0,0.75) circle (1); \fill[cyan!18] (-0.65,-0.4) circle (1);\end{scope}
    \draw (0,0.75) circle (1); \draw (-0.65,-0.4) circle (1); \draw (0.65,-0.4) circle (1);
    \node at (0,1.95) {$A$}; \node at (-1.5,-1.35) {$B$}; \node at (1.5,-1.35) {$C$};
    \node at (0,-2.05) {$(A\cap B)$};
  \end{scope}
  \node at (2.25,-2.05) {$\cap$};
  \begin{scope}[shift={(4.5,0)}]
    \fill[cyan!18] (0.65,-0.4) circle (1);
    \draw (0,0.75) circle (1); \draw (-0.65,-0.4) circle (1); \draw (0.65,-0.4) circle (1);
    \node at (0,1.95) {$A$}; \node at (-1.5,-1.35) {$B$}; \node at (1.5,-1.35) {$C$};
    \node at (0,-2.05) {$C$};
  \end{scope}
  \begin{scope}[shift={(9.0,0)}]
    \draw (0,0.75) circle (1); \draw (-0.65,-0.4) circle (1); \draw (0.65,-0.4) circle (1);
    \node at (0,1.95) {$A$}; \node at (-1.5,-1.35) {$B$}; \node at (1.5,-1.35) {$C$};
    \node at (0,-2.05) {$=(A\cap B)\cap C$};
  \end{scope}
  \node[anchor=west] at (-3.7,-2.9) {[우변]};
  \begin{scope}[shift={(0,-4.6)}]
    \fill[cyan!18] (0,0.75) circle (1);
    \draw (0,0.75) circle (1); \draw (-0.65,-0.4) circle (1); \draw (0.65,-0.4) circle (1);
    \node at (0,1.95) {$A$}; \node at (-1.5,-1.35) {$B$}; \node at (1.5,-1.35) {$C$};
    \node at (0,-2.05) {$A$};
  \end{scope}
  \node at (2.25,-6.65) {$\cap$};
  \begin{scope}[shift={(4.5,-4.6)}]
    \begin{scope}\clip (-0.65,-0.4) circle (1); \fill[cyan!18] (0.65,-0.4) circle (1);\end{scope}
    \draw (0,0.75) circle (1); \draw (-0.65,-0.4) circle (1); \draw (0.65,-0.4) circle (1);
    \node at (0,1.95) {$A$}; \node at (-1.5,-1.35) {$B$}; \node at (1.5,-1.35) {$C$};
    \node at (0,-2.05) {$(B\cap C)$};
  \end{scope}
  \begin{scope}[shift={(9.0,-4.6)}]
    \draw (0,0.75) circle (1); \draw (-0.65,-0.4) circle (1); \draw (0.65,-0.4) circle (1);
    \node at (0,1.95) {$A$}; \node at (-1.5,-1.35) {$B$}; \node at (1.5,-1.35) {$C$};
    \node at (0,-2.05) {$=A\cap(B\cap C)$};
  \end{scope}
\end{tikzpicture}
